\section*{Orthogonal elements:}

Let \(\Phi\) be a form \(\varepsilon\)-hermitian on a left A-module E. One says that two elements \(x, y\) of E are orthogonal (for \(\Phi\) ) if \(\Phi(x, y) = 0\) ; this relation is symmetric in \(x\) and \(y\) . For every submodule M of E, the set of \(x \in E\) which are orthogonal to all elements of M is a submodule denoted \(\mathbf{M}^{\circ}\) and called the orthogonal of the submodule M. One says that \(\Phi\) is non-degenerate if \(\mathbf{E}^{\circ} = \{0\}\) . When E is a finite-dimensional vector space and \(\Phi\) is non-degenerate, one has codim \(\mathbf{M}^{\circ} = \dim \mathbf{M}\) and \(\mathbf{M}^{\circ\circ} = \mathbf{M}\) for every subspace M of E; moreover, for every pair of subspaces M, N of E, one has

\[
(\mathbf {M} + \mathbf {N}) ^ {0} = \mathbf {M} ^ {0} \cap \mathbf {N} ^ {0}, \quad (\mathbf {M} \cap \mathbf {N}) ^ {0} = \mathbf {M} ^ {0} + \mathbf {N} ^ {0}.
\]

Suppose the ring A is commutative, and let Q be a quadratic form on the A-module E, \(\Phi\) the bilinear form associated with Q. One says that two elements of E are orthogonal (for Q) if they are orthogonal for \(\Phi\) ; the orthogonal of a submodule M (for Q) is the orthogonal M \(^{o}\) of M for \(\Phi\) . One says that Q is non-degenerate if \(\Phi\) is non-degenerate.

